% ===== English content body =====
\begin{titlepage}\centering
\vspace*{4mm}
{\displayfont\bfseries\fontsize{40}{44}\selectfont Balkon\,Borg}\\[3mm]
{\displayfont\LARGE Assembly manual — hardware}\\[1mm]
{\large\color{soft} from the parts to the finished device}\\[9mm]
\dekocover{cover.png}{Cover: the finished device under the balcony}\\[7mm]
\includegraphics[width=0.6\textwidth]{enclosure.png}\\[8mm]
\begin{tcolorbox}[width=0.86\textwidth,colback=cfun,colframe=ink,arc=4pt]\centering
This manual goes in order: solder, wire up, install, mount. None of it is hard if you take
it step by step. Take your time and work through the chapters in sequence.
\end{tcolorbox}
\vfill
{\footnotesize\color{soft}Technical images from the CAD/PCB of the \texttt{balkon-borg} repo
(not drawn) · as of \today}
\end{titlepage}

\section*{Before you start}
This manual covers the full \textbf{hardware build}: populate the board, make up the
cables, everything into the enclosure, wire it, mount it on the ceiling. The software
(flashing the ESP, setting up the Pi) is deliberately not included. Set aside a calm
afternoon, with a break between soldering and wiring.

\begin{achtung}
The \textbf{230\,V power supply stays outside} in its own housing. Only \textbf{5\,V} goes
into the printed enclosure. Work on the 230\,V part only when it is disconnected from mains.
\end{achtung}
\begin{warum}
Plastic (PA12) and mains voltage do not mix: no defined fire rating, no creepage distances.
So the PSU sits separately and the printed enclosure carries low voltage only. The enclosure
comes out of the printer in \textbf{two halves} (seam in the middle); during the build two
4\,mm dowel pins align them and three internal M3 clamps bolt them together.
\end{warum}

\tableofcontents

% ---------------- Parts list ----------------
\clearpage
\section{Parts list}
\dekopikto{p-tools.png}{parts}
\begin{ziel}
Have everything complete and pre-sorted before you start. Quantities are per \emph{one}
device. Authoritative reference: \texttt{pcb/docs/board-spec.md} (board) and
\texttt{docs/cost-estimate.md} (whole device).
\end{ziel}

\textbf{On the board (to solder):}\par
\rowcolors{2}{rowb}{white}
\begin{tabular}{clp{88mm}}
\toprule
\textbf{\#} & \textbf{Part} & \textbf{Description} \\ \midrule
1 & ESP32-DevKitC-V4 & official Espressif board, WROOM-32E, 38-pin, \textbf{25.4\,mm} row spacing (not a 0.9" clone) \\
2 & female header 1×19 & 2.54\,mm, for the ESP (socketed) \\
4 & transistor BC337 & NPN, TO-92, drives the button LEDs \\
2 & resistor 220\,$\Omega$ & series resistor, radar UART \\
7 & resistor 100\,$\Omega$ & series resistor, encoder + buttons \\
4 & resistor 1\,k$\Omega$ & transistor base resistor \\
2 & resistor 4.7\,k$\Omega$ & I²C pull-ups (optional / DNP, often on the BME breakout) \\
1 & electrolytic 10\,\textmu F & 3V3 decoupling (\textbf{polarised}) \\
1 & capacitor 100\,nF & 3V3 decoupling (non-polarised) \\
1 & polyfuse ~2\,A & resettable 5\,V input fuse \\
1 & JST-XH header 2p & \texttt{J\_PWR} (5\,V input) \\
7 & JST-XH header 4p & \texttt{J\_RADAR}, \texttt{J\_BME}, \texttt{J\_ENC}, \texttt{J\_BTN1--4} \\
\bottomrule
\end{tabular}

\textbf{External modules:}\par
\rowcolors{2}{rowb}{white}
\begin{tabular}{clp{88mm}}
\toprule
\textbf{\#} & \textbf{Module} & \textbf{Description} \\ \midrule
1 & Raspberry Pi 5 (8\,GB) & + Active Cooler + microSD — the central compute unit \\
1 & LD2410B & 24-GHz presence radar (UART), also senses a still-seated person \\
1 & BME280 & genuine Bosch sensor: temperature / humidity / pressure (I²C) \\
1 & EC11 encoder & rotary encoder with push-button (brightness) \\
4 & 12\,mm button & illuminated, stainless, \textbf{5\,V} ring LED \\
1 & WLED controller & Athom high-power (pre-flashed) \\
1 & SK6812 RGBW-WW & panel 8×43 = 344 px, on an aluminium plate + opal diffuser \\
1 & RTL-SDR V3 & + telescopic SMA antenna (ADS-B, optional LoRa) \\
1 & USB microphone & on the Pi (BirdNET) \\
1 & audio chain & USB sound card (DELOCK 61645) + PAM8403 amp + Visaton BF\,45 \\
1 & Camera Module 3 & on the Pi's CSI port \\
1 & Mean Well LRS-150-5 & 5\,V/22\,A PSU, \textbf{external} in its own housing \\
\bottomrule
\end{tabular}

\textbf{Small parts / mechanical:}\par
\rowcolors{2}{rowb}{white}
\begin{tabular}{clp{88mm}}
\toprule
\textbf{\#} & \textbf{Part} & \textbf{Description} \\ \midrule
~8 & threaded insert M2.5 & heat-set (Pi, carrier board, WLED holder) \\
3 & threaded insert M3 & heat-set, seam clamps joining the two halves \\
2 & dowel pin 4\,mm & aligns the halves at the seam \\
-- & screws M2.5 / M3 / M5 & M2.5 inside, M3 ×3 seam clamps, M5 ×4 ceiling ears \\
1 & XT60 socket & 5\,V feed-in in the rear wall \\
2 & Wago 221 (5-way) & 5\,V+ and GND star distribution \\
-- & fuses + holders & inline 10\,A (LED), 5\,A (Pi) \\
-- & wire + ferrules & 2.5 / 1.5 / 0.5\,mm² \\
-- & JST-XH crimp set & contacts + housings for the sensor/control cables \\
-- & misc & panel nuts, thermal pad, glue, insect mesh, zip ties \\
\bottomrule
\end{tabular}

\begin{schmunzel}
Sort the resistors by value into little dishes right away. Telling four values apart on the
tape later costs more nerves than the rest of the build.
\end{schmunzel}

% ---------------- Tools ----------------
\clearpage
\section{Tools}
\dekopikto{p-solder.png}{tools}
\begin{itemize}[leftmargin=8mm,itemsep=2pt]
  \item \faBolt\ temperature-controlled soldering iron (fine tip) + solder, desoldering braid
  \item \faCut\ side cutter, wire stripper, \textbf{JST-XH crimping tool}
  \item \faScrewdriver\ small Phillips screwdriver, tweezers
  \item \faFire\ soldering iron or hot air for the \textbf{threaded inserts}
  \item \faSearch\ multimeter (continuity) for the final visual check
\end{itemize}
\begin{tipp}
A 30-euro temperature-controlled iron is plenty (~330\,°C for lead-free solder). The JST
crimping tool is the only special tool and saves the most time.
\end{tipp}

% ---------------- A Board ----------------
\clearpage
\section{Populate the board (soldering)}
\dekopikto{p-solder.png}{solder}
\begin{ziel}
The empty board becomes the fully populated carrier with the ESP32 seated on it. Populate
\textbf{low to tall}, so everything sits flat whenever you flip the board over.
\end{ziel}
\begin{schritte}
  \item \textbf{15 resistors} (R1–R15): insert all, bend the legs slightly, then solder in
        one pass and trim the ends. Resistors are non-polarised — either way round.
  \item \textbf{4 transistors} BC337 (Q1–Q4) — \emph{mind the orientation} (see caution box).
  \item \textbf{Capacitors}: C1 (10\,\textmu F electrolytic, \textbf{polarised} — the minus
        stripe goes to the GND side) and C2 (100\,nF, non-polarised).
  \item \textbf{Polyfuse F1} (the 5\,V input fuse).
  \item \textbf{8 JST-XH headers}: \texttt{J\_PWR} (2p) plus \texttt{J\_RADAR},
        \texttt{J\_BME}, \texttt{J\_ENC}, \texttt{J\_BTN1--4} (4p each).
  \item \textbf{2 female headers} (1×19) for the ESP — align them straight before soldering.
  \item \textbf{Seat the ESP32-DevKitC-V4}.
\end{schritte}
\bild{pcb-top.png}{0.82\textwidth}{The populated board with part designators (source: KiCad render).}
\begin{achtung}
BC337: the \textbf{flat face} of the transistor must face the silkscreen marking on the
board. The wrong way round, the button LEDs never switch. Look first, then solder. And put
the electrolytic C1 in the right way — the minus stripe points away from the 3V3 marking.
\end{achtung}
\begin{tipp}
Female headers like to tilt. Seat the ESP on them \emph{before} soldering, square everything
up, then solder just one outer pin each and correct before you solder the rest.
\end{tipp}
\begin{ergebnis}
A fully populated, DRC-clean carrier board with the ESP32 seated.
\end{ergebnis}

% ---------------- B Cables ----------------
\clearpage
\section{Make up the cables (crimping)}
\dekopikto{p-crimp.png}{crimp}
\begin{ziel}
All connecting cables from the board to sensors, controls and power — with the right pinout
and wire colour. \textbf{Note:} the JST cables are \emph{crimped}, not soldered. You only
solder where a module has bare pads instead of a connector.
\end{ziel}
\begin{schritte}
  \item Measure cable lengths roughly (generous; there is room in the enclosure).
  \item Strip ~3\,mm and \textbf{crimp the JST-XH contacts} (crimping tool).
  \item Push the contacts into the housing until they click home.
  \item Keep wire colours consistent: \ader{wRD}{5\,V} \quad \ader{wBK}{GND} \quad
        \ader{wOG}{3.3\,V} \quad \ader{wGN}{signal} \quad \ader{wYE}{signal} \quad
        \ader{wBU}{signal} \quad \ader{wWH}{signal}.
\end{schritte}
\bild{wiring-harness.png}{0.78\textwidth}{The full cable harness: which wire to which pin, with colours (source: netlist → WireViz).}
\begin{warum}
Consistent wire colours (red = 5\,V, black = GND) are not decoration: they stop you from
swapping plus and minus when you plug things together — the most common and most annoying
mistake, and hardly possible with a colour code.
\end{warum}

% ---------------- C Sensors + tables ----------------
\clearpage
\section{Connect sensors \& controls}
\dekopikto{p-connect.png}{connect}
\begin{ziel}
Every cable gets its module — wire by wire. Below, each connector lists exactly which pin
goes over which wire colour to which module pin.
\end{ziel}

\bild{board-annotated.png}{0.86\textwidth}{Which connector goes to which module (board render with callouts).}

\textbf{Example as a diagram — radar (\texttt{J\_RADAR} → LD2410B):}
\begin{center}\begin{tikzpicture}
  \node[wpin,fill=ink,text=white] at (0,0.6) {J\_RADAR};
  \node[wpin,fill=ink,text=white] at (7,0.6) {LD2410B};
  \wire{1}{1 · 5V}{wRD}{VCC}
  \wire{2}{2 · GND}{wBK}{GND}
  \wire{3}{3 · RX-sig.}{wGN}{RX}
  \wire{4}{4 · TX-sig.}{wYE}{TX}
\end{tikzpicture}\end{center}

\textbf{All connector pinouts (board pin — wire — module pin):}\par
\rowcolors{2}{rowb}{white}
\begin{tabular}{llll}
\toprule
\textbf{Connector} & \textbf{Board pin} & \textbf{Wire} & \textbf{Module pin} \\ \midrule
\texttt{J\_RADAR} & 1 · 5V & \ader{wRD}{} & LD2410B VCC \\
 & 2 · GND & \ader{wBK}{} & GND \\
 & 3 · RX-signal & \ader{wGN}{} & RX \\
 & 4 · TX-signal & \ader{wYE}{} & TX \\
\texttt{J\_BME} & 1 · 3V3 & \ader{wOG}{} & BME280 VCC \\
 & 2 · GND & \ader{wBK}{} & GND \\
 & 3 · SCL & \ader{wWH}{} & SCL \\
 & 4 · SDA & \ader{wGN}{} & SDA \\
\texttt{J\_ENC} & 1 · A & \ader{wGN}{} & EC11 A \\
 & 2 · B & \ader{wYE}{} & B \\
 & 3 · SW & \ader{wBU}{} & push-button \\
 & 4 · GND & \ader{wBK}{} & GND \\
\texttt{J\_BTN1--4} & 1 · SW & \ader{wGN}{} & button switch \\
 & 2 · GND & \ader{wBK}{} & switch COM \\
 & 3 · 5V & \ader{wRD}{} & LED\,+ (anode) \\
 & 4 · LED-K & \ader{wWH}{} & LED\,-- (cathode) \\
\texttt{J\_PWR} & 1 · 5V & \ader{wRD}{} & PSU 5\,V (Wago) \\
 & 2 · GND & \ader{wBK}{} & PSU GND (Wago) \\
\bottomrule
\end{tabular}
\begin{tipp}
Buy the \textbf{5\,V variant} of the button ring LED — then the driver on the board needs
no extra resistor. If a module has bare solder pads instead of a connector (some radar/BME
breakouts), solder the wire there directly — same order and colour as in the table.
\end{tipp}
\begin{ergebnis}
All four sensor/control modules are correctly wired to their cables.
\end{ergebnis}

% ---------------- D Inserts ----------------
\clearpage
\section{Prepare the enclosure: set inserts, join the halves}
\dekopikto{p-insert.png}{set inserts}
\begin{ziel}
The threaded inserts sit, the two enclosure halves are bolted together — and the LED panel
already lies inside, because after joining it no longer fits through any opening.
\end{ziel}
\begin{schritte}
  \item Set the \textbf{M2.5 inserts} for Pi 5, carrier board and WLED holder (~8 of them).
  \item \textbf{M3 inserts} into the three seam clamps of one half (2× rear wall, 1× front
        of the floor). The \textbf{ceiling ears} have M5 through-holes — no insert needed.
  \item \textbf{4\,mm dowel pins} into the dowel posts of one half (a drop of glue holds them).
  \item \textbf{Lay in the LED panel, then join:} slide the panel (with its aluminium
        plate) in centred while bringing the halves together, holding it about 2\,cm
        behind the front edge (there it clears both the front seam clamp and the locating
        nubs). Push the halves onto the pins and bolt them with the three \textbf{M3
        screws} — from the open front, long screwdriver.
  \item Then slide the panel \textbf{forward} between the four locating nubs (it is glued
        later, together with the diffuser).
\end{schritte}
\bild{half-left.png}{0.58\textwidth}{View of one enclosure half (interior visible).}
\begin{achtung}
Sink the inserts with the soldering iron \textbf{straight and not too hot}. Set crooked,
the thread pulls out later — slow and vertical is better.
\end{achtung}
\begin{ergebnis}
Every mounting point has a solid M2.5 thread.
\end{ergebnis}

% ---------------- E Install ----------------
\clearpage
\section{Install the components}
\dekopikto{p-install.png}{install}
\begin{ziel}
Every module sits in its place — on ~8\,mm standoffs, with plenty of room to the LED panel.
\end{ziel}
\begin{schritte}
  \item Screw the \textbf{Raspberry Pi 5} onto the rear-wall bosses.
  \item \textbf{Carrier board} into the right-hand bay on its M2.5 bosses (parts toward the front).
  \item Glue the \textbf{LED panel} (inside since the joining step, on its aluminium plate
        as heatsink) to the frame back between the nubs. Glue the \textbf{opal diffuser}
        into the rebate only as the very last step — after that the interior is reachable
        only through the ceiling opening.
  \item \textbf{Tower LEDs + radar} into the downward tower; glue the radar over the
        30×5.5 slot from inside (a bead all round — holds it and seals against insects).
  \item \textbf{Camera} into the hanging box: lens block flat against the tilted front
        wall, barrel through the bore, tack the board corners with \textbf{hot glue}. Plus
        \textbf{microphone} and \textbf{audio} (USB sound card, PAM8403 amp, Visaton speaker).
  \item \textbf{RTL-SDR} on a \emph{USB-2} port, antenna to the SMA socket.
\end{schritte}
\begin{figure}[htbp]\centering
  \includegraphics[width=0.47\textwidth]{view-inside.png}\hfill
  \includegraphics[width=0.47\textwidth]{view-bottom.png}\par
  {\footnotesize\color{soft}Interior (left) and underside with camera/tower opening (right).}
\end{figure}
\begin{warum}
The SDR belongs on USB-2, not USB-3: USB-3 radiates right on 1090\,MHz and would jam the
aircraft reception (ADS-B). Small move, big difference.
\end{warum}
\begin{ergebnis}
A fully populated enclosure — just not wired yet.
\end{ergebnis}

% ---------------- F Wiring ----------------
\clearpage
\section{Power \& wiring}
\dekopikto{p-wire.png}{wire up}
\begin{ziel}
The 5\,V star distribution is in place (with fuses), and every cable from the board to the
sensors is plugged in.
\end{ziel}
\begin{schritte}
  \item \textbf{XT60 socket} into the rear wall, 5\,V from the external PSU comes here.
  \item \textbf{Wago blocks} for 5\,V+ and GND (common ground for everything).
  \item \textbf{Fuses} inline: 10\,A (LED panel), 5\,A (Pi\,5), 2\,A (board, via its F1).
  \item Branches with the right gauge: 2.5 / 1.5 / 0.5\,mm².
  \item \textbf{Audio amp} tap 5\,V/GND off the \emph{Pi branch} (same reference → no hum).
  \item Plug all \textbf{JST cables} board\,↔\,sensors — pinout per the table in chapter~5.
\end{schritte}
\begin{achtung}
Before first power-up: \textbf{check polarity}. Applying 5\,V reversed kills the Pi and the
board instantly. Red = 5\,V, black = GND, everywhere.
\end{achtung}
\begin{tipp}
Keep the 5\,V line short and thick, especially to the LED panel — otherwise brightness drops
at the far end of the strip. The long run belongs on the 230\,V side (to the external PSU).
\end{tipp}
\begin{ergebnis}
A wired device with fused power and common ground.
\end{ergebnis}

% ---------------- G Mounting ----------------
\clearpage
\section{Ceiling mounting}
\dekopikto{p-assemble.png}{mount}
\begin{ziel}
The enclosure hangs securely under the balcony. The halves have been bolted together since
the preparation step — all that is left is the ceiling.
\end{ziel}
\begin{schritte}
  \item Final \textbf{cable check}: nothing pinched, nothing kinked.
  \item Screw the device to the ceiling from below through the \textbf{ceiling ears}
        (M5, ×4) — a washer under each head helps.
\end{schritte}
\begin{warum}
The four ears carry the entire weight. Use anchors matched to the ceiling material; when in
doubt, go one size up.
\end{warum}
\begin{ergebnis}
A closed, mounted Balkon\,Borg in its place.
\end{ergebnis}

% ---------------- H Check ----------------
\clearpage
\section{Visual check before power-up}
\dekopikto{p-check.png}{check}
\begin{ziel}
Go over everything once before power is applied. The software is a separate chapter — this
is hardware only.
\end{ziel}
\begin{schritte}
  \item Polarity of all 5\,V connections checked (red/black)?
  \item No solder bridges/shorts on the board (continuity tester)?
  \item All JST connectors fully clicked home, pinout as in chapter~5?
  \item Fuses (10/5\,A) in place?
  \item The 230\,V PSU is outside and still \textbf{disconnected from mains}?
\end{schritte}
\begin{ergebnis}
A built, checked device — ready for the software (separate manual).
\end{ergebnis}
